%% ============================================================
%%  Chapter 5 — Release 2: Learning Experience and Course Access
%% ============================================================
\begingroup
\renewcommand{\baselinestretch}{0.85}
\chapter{Release 2: Learning Experience and Course Access}
\minitoc
\endgroup
\newpage

%% ------------------------------------------------------------
\section*{Introduction}
\addcontentsline{toc}{section}{Introduction}
%% ------------------------------------------------------------

With the foundational infrastructure established in Release 1, Release 2 shifts the
focus toward the student-facing dimension of the platform. While the first release
provided the tools necessary to create and structure educational content, this release
delivers the mechanisms through which students can discover, access, and progress
through that content. It represents the core of the learning experience that Tesla
Academy is designed to offer.

Release 2 corresponds to Sprints 3 and 4 of the project plan. Sprint 3 introduces the
assignment workflow, lesson access and completion tracking, progress measurement, and
student notifications. Sprint 4 extends the platform with certificate generation,
role-specific dashboards for both students and instructors, the public course catalog
with search and filtering capabilities, course reviews, and free course enrollment.

The chapter follows the same structure as the previous one: sprint backlog, functional
and technical specifications for all user stories, design models, and implementation
results.

%% ------------------------------------------------------------
\section{Sprint Backlog}
%% ------------------------------------------------------------

\begin{table}[H]
\centering
\caption{Sprint Backlog — Release 2}
\label{tab:backlog_r2}
\renewcommand{\arraystretch}{1.25}
\footnotesize
\begin{tabularx}{\textwidth}{|p{1.4cm}|p{1.5cm}|p{2.8cm}|p{6.8cm}|p{1.5cm}|}
\hline
\textbf{Sprint} & \textbf{Code} & \textbf{Feature} & \textbf{User Story} & \textbf{Est. (h)} \\
\hline
\multirow{12}{*}{Sprint 3}
  & US-019 & Assignment creation & Create an assignment with title, instructions, and deadline & 16 \\
  \cline{2-5}
  & US-020 & Assignment resources & Attach files and resources to an assignment & 8 \\
  \cline{2-5}
  & US-021 & Assignment publication & Publish an assignment to enrolled students & 8 \\
  \cline{2-5}
  & US-022 & Assignment submission & Submit an assignment response before the deadline & 16 \\
  \cline{2-5}
  & US-023 & Assignment review & Review submissions and provide feedback and grades & 16 \\
  \cline{2-5}
  & US-024 & My courses & View the list of enrolled courses & 8 \\
  \cline{2-5}
  & US-025 & Access lesson & Access lesson content such as videos and documents & 16 \\
  \cline{2-5}
  & US-026 & Complete lesson & Mark a lesson as completed & 8 \\
  \cline{2-5}
  & US-027 & Progress tracking & Track course progress percentage & 8 \\
  \cline{2-5}
  & US-028 & Consult assignments & View the list of assignments for enrolled courses & 8 \\
  \cline{2-5}
  & US-029 & Assignment feedback & Consult grades and feedback for submitted assignments & 8 \\
  \cline{2-5}
  & US-030 & Announcements & Consult announcements published by the instructor & 8 \\
\hline
\multirow{11}{*}{Sprint 4}
  & US-031 & Certificate generation & Obtain a certificate after completing a course & 16 \\
  \cline{2-5}
  & US-032 & Certificate download & Download the certificate in PDF format & 8 \\
  \cline{2-5}
  & US-033 & Student dashboard & View courses, progress, assignments, and certificates & 16 \\
  \cline{2-5}
  & US-034 & Instructor dashboard & Monitor course statistics and student submissions & 16 \\
  \cline{2-5}
  & US-035 & Publish announcement & Publish announcements for enrolled students & 8 \\
  \cline{2-5}
  & US-036 & Course catalog & View the list of available courses & 16 \\
  \cline{2-5}
  & US-037 & Course search & Search for a course by keyword & 8 \\
  \cline{2-5}
  & US-038 & Catalog filters & Filter courses by category, level, or price & 8 \\
  \cline{2-5}
  & US-039 & Course details & View detailed information about a course & 8 \\
  \cline{2-5}
  & US-040 & Reviews and ratings & Rate and comment on a completed course & 8 \\
  \cline{2-5}
  & US-041 & Free course enrollment & Enroll in a free course without payment & 8 \\
\hline
\end{tabularx}
\end{table}

%% ------------------------------------------------------------
\section{Functional and Technical Specifications}
%% ------------------------------------------------------------

This section provides a detailed breakdown of all user stories implemented in Release 2.
For each functional area, a specification table describes the user story, the associated
business rules, the technical implementation details, and the acceptance criteria that
were used to validate the delivered functionality.

%% ----------------------------
\subsection{Assignment Workflow (US-019 to US-023)}
%% ----------------------------

The assignment module introduces a structured assessment workflow involving three
actors: the instructor who creates and publishes assessments, the student who submits
responses, and the system which handles deadline enforcement and feedback notification.
The workflow supports multiple question types and accommodates both text-based and
file-based submission formats, providing instructors with flexible evaluation tools.

\begin{table}[H]
\centering
\caption{Specification — US-019: Assignment Creation}
\label{tab:us019}
\renewcommand{\arraystretch}{1.2}
\begin{tabularx}{\textwidth}{|p{3.5cm}|X|}
\hline
\textbf{User Story} & As an Instructor, I want to create an assignment with a title, instructions, and a deadline so that students have a clear and time-bound assessment task. \\
\hline
\textbf{Precondition} & The instructor is authenticated and owns at least one published course. \\
\hline
\textbf{Business Rules} &
\begin{itemize}[noitemsep, topsep=2pt]
  \item An assignment must be associated with an existing course;
  \item The assignment must include a title, instructions, and a deadline;
  \item Assignments are created in an unpublished state and are not visible to students until explicitly published.
\end{itemize} \\
\hline
\textbf{Technical Specification} &
\begin{itemize}[noitemsep, topsep=2pt]
  \item Assignments are stored in the \texttt{Assignment} model, linked to a \texttt{Batch} and scoped by \texttt{organizationId};
  \item Questions are stored in the \texttt{AssignmentQuestion} model and support four types: \texttt{QUIZ}, \texttt{TRUE\_FALSE}, \texttt{SHORT\_ANSWER}, and \texttt{FILE\_SUBMISSION};
  \item The deadline is stored as a timestamp and enforced by the backend at submission time.
\end{itemize} \\
\hline
\textbf{Acceptance Criteria} &
\begin{itemize}[noitemsep, topsep=2pt]
  \item The instructor can create an assignment with all required fields;
  \item The assignment is saved in unpublished state and not visible to students;
  \item Missing required fields are rejected with appropriate validation messages.
\end{itemize} \\
\hline
\end{tabularx}
\end{table}

\begin{table}[H]
\centering
\caption{Specification — US-020: Assignment Resources}
\label{tab:us020}
\renewcommand{\arraystretch}{1.2}
\begin{tabularx}{\textwidth}{|p{3.5cm}|X|}
\hline
\textbf{User Story} & As an Instructor, I want to attach files and supporting resources to an assignment so that students have all necessary material to complete their work. \\
\hline
\textbf{Precondition} & The assignment has been created and the instructor is authenticated. \\
\hline
\textbf{Business Rules} &
\begin{itemize}[noitemsep, topsep=2pt]
  \item Resources must be attached to an existing assignment;
  \item Supported resource types include documents and reference files;
  \item Resources remain accessible to enrolled students once the assignment is published.
\end{itemize} \\
\hline
\textbf{Technical Specification} &
\begin{itemize}[noitemsep, topsep=2pt]
  \item Attached files are uploaded to Supabase Storage via presigned URLs;
  \item File URLs are stored in the \texttt{Assignment} model as a JSON array;
  \item The backend validates ownership before accepting resource attachments.
\end{itemize} \\
\hline
\textbf{Acceptance Criteria} &
\begin{itemize}[noitemsep, topsep=2pt]
  \item The instructor can attach one or more files to an assignment;
  \item Attached resources are stored successfully and linked to the assignment;
  \item Students can access the resources once the assignment is published.
\end{itemize} \\
\hline
\end{tabularx}
\end{table}

\begin{table}[H]
\centering
\caption{Specification — US-021: Assignment Publication}
\label{tab:us021}
\renewcommand{\arraystretch}{1.2}
\begin{tabularx}{\textwidth}{|p{3.5cm}|X|}
\hline
\textbf{User Story} & As an Instructor, I want to publish an assignment so that enrolled students can view it and submit their responses. \\
\hline
\textbf{Precondition} & The assignment exists and contains at least one question. \\
\hline
\textbf{Business Rules} &
\begin{itemize}[noitemsep, topsep=2pt]
  \item An assignment must be explicitly published before students can access it;
  \item Publication is independent of the course publication state;
  \item A published assignment cannot be deleted if students have already submitted responses.
\end{itemize} \\
\hline
\textbf{Technical Specification} &
\begin{itemize}[noitemsep, topsep=2pt]
  \item The publication state is managed through a boolean field on the \texttt{Assignment} model;
  \item Student-facing endpoints filter assignments by publication status before returning results;
  \item The backend validates that the assignment contains at least one question before allowing publication.
\end{itemize} \\
\hline
\textbf{Acceptance Criteria} &
\begin{itemize}[noitemsep, topsep=2pt]
  \item The instructor can publish an assignment from the management interface;
  \item Published assignments appear in the student workspace;
  \item Unpublished assignments remain hidden from students.
\end{itemize} \\
\hline
\end{tabularx}
\end{table}

\begin{table}[H]
\centering
\caption{Specification — US-022: Assignment Submission}
\label{tab:us022}
\renewcommand{\arraystretch}{1.2}
\begin{tabularx}{\textwidth}{|p{3.5cm}|X|}
\hline
\textbf{User Story} & As a Student, I want to submit my assignment response with text or attached files before the deadline so that the instructor can evaluate my work. \\
\hline
\textbf{Precondition} & The student is enrolled in the course and the assignment is published with an active deadline. \\
\hline
\textbf{Business Rules} &
\begin{itemize}[noitemsep, topsep=2pt]
  \item Students can only submit if they are enrolled in the associated course;
  \item Submissions are rejected after the deadline has passed;
  \item A student can only submit one response per assignment.
\end{itemize} \\
\hline
\textbf{Technical Specification} &
\begin{itemize}[noitemsep, topsep=2pt]
  \item Submissions are stored in the \texttt{AssignmentSubmission} table;
  \item The \texttt{fileUrls} JSON field stores the URLs of uploaded attachment files;
  \item The backend validates enrollment and deadline before persisting the submission;
  \item File uploads follow the same presigned URL mechanism used elsewhere on the platform.
\end{itemize} \\
\hline
\textbf{Acceptance Criteria} &
\begin{itemize}[noitemsep, topsep=2pt]
  \item An enrolled student can submit a response before the deadline;
  \item Submissions after the deadline are rejected with an appropriate error;
  \item The submission is stored and visible to the instructor for review.
\end{itemize} \\
\hline
\end{tabularx}
\end{table}

\begin{table}[H]
\centering
\caption{Specification — US-023: Assignment Review}
\label{tab:us023}
\renewcommand{\arraystretch}{1.2}
\begin{tabularx}{\textwidth}{|p{3.5cm}|X|}
\hline
\textbf{User Story} & As an Instructor, I want to review student submissions and provide feedback and grades so that students can understand their performance. \\
\hline
\textbf{Precondition} & At least one student has submitted a response to the assignment. \\
\hline
\textbf{Business Rules} &
\begin{itemize}[noitemsep, topsep=2pt]
  \item Instructors can review all submissions for assignments within their courses;
  \item Each submission can receive a numeric grade and written feedback;
  \item A notification email is sent to the student upon grading.
\end{itemize} \\
\hline
\textbf{Technical Specification} &
\begin{itemize}[noitemsep, topsep=2pt]
  \item Grades and feedback are stored in the \texttt{marksAwarded} and \texttt{feedback} fields of the \texttt{AssignmentSubmission} record;
  \item An email notification is dispatched to the student upon saving the grade;
  \item The review interface is accessible only to instructors owning the associated course.
\end{itemize} \\
\hline
\textbf{Acceptance Criteria} &
\begin{itemize}[noitemsep, topsep=2pt]
  \item The instructor can view all submissions for an assignment;
  \item The instructor can assign a grade and written feedback to each submission;
  \item The student receives an email notification after grading.
\end{itemize} \\
\hline
\end{tabularx}
\end{table}

%% ----------------------------
\subsection{Learning Experience (US-024 to US-030)}
%% ----------------------------

The learning experience module covers all interactions a student has with enrolled
courses. It includes accessing the list of enrolled courses, navigating lesson content,
marking topics as complete, tracking progress, consulting assignment tasks and feedback,
and reading instructor announcements. These features collectively form the daily
learning workflow of a student on the platform.

\begin{table}[H]
\centering
\caption{Specification — US-024: My Courses}
\label{tab:us024}
\renewcommand{\arraystretch}{1.2}
\begin{tabularx}{\textwidth}{|p{3.5cm}|X|}
\hline
\textbf{User Story} & As a Student, I want to view the list of courses I am enrolled in so that I can easily navigate to my learning content. \\
\hline
\textbf{Precondition} & The student is authenticated and is enrolled in at least one course. \\
\hline
\textbf{Business Rules} &
\begin{itemize}[noitemsep, topsep=2pt]
  \item Only courses for which the student holds an active enrollment record are displayed;
  \item The list shows basic course information including title, progress, and instructor name.
\end{itemize} \\
\hline
\textbf{Technical Specification} &
\begin{itemize}[noitemsep, topsep=2pt]
  \item The backend queries \texttt{BatchEnrollment} records filtered by the authenticated student's identifier;
  \item Course progress is computed server-side and included in the response;
  \item Results are scoped by \texttt{organizationId} to enforce tenant isolation.
\end{itemize} \\
\hline
\textbf{Acceptance Criteria} &
\begin{itemize}[noitemsep, topsep=2pt]
  \item The student can view all their enrolled courses in a dedicated workspace section;
  \item Each course card displays the title, instructor, and current progress;
  \item Courses from other organizations are never displayed.
\end{itemize} \\
\hline
\end{tabularx}
\end{table}

\begin{table}[H]
\centering
\caption{Specification — US-025: Access Lesson}
\label{tab:us025}
\renewcommand{\arraystretch}{1.2}
\begin{tabularx}{\textwidth}{|p{3.5cm}|X|}
\hline
\textbf{User Story} & As a Student, I want to access lesson content such as videos, documents, and text so that I can follow the course material at my own pace. \\
\hline
\textbf{Precondition} & The student is enrolled in the course and the lesson belongs to a published course. \\
\hline
\textbf{Business Rules} &
\begin{itemize}[noitemsep, topsep=2pt]
  \item A student must be enrolled in a course to access any of its lesson content;
  \item Content types include uploaded videos, YouTube videos, PDF documents, rich-text blocks, and external links;
  \item Access to content files is controlled through the backend to prevent unauthorized downloads.
\end{itemize} \\
\hline
\textbf{Technical Specification} &
\begin{itemize}[noitemsep, topsep=2pt]
  \item The backend verifies enrollment before returning content URLs or data;
  \item Video files are served through Supabase Storage URLs;
  \item The lesson viewer renders each content type using its corresponding component on the frontend.
\end{itemize} \\
\hline
\textbf{Acceptance Criteria} &
\begin{itemize}[noitemsep, topsep=2pt]
  \item An enrolled student can access all lesson content within a published course;
  \item Non-enrolled users are denied access and redirected appropriately;
  \item All content types are rendered correctly in the lesson viewer.
\end{itemize} \\
\hline
\end{tabularx}
\end{table}

\begin{table}[H]
\centering
\caption{Specification — US-026 and US-027: Lesson Completion and Progress Tracking}
\label{tab:us026}
\renewcommand{\arraystretch}{1.2}
\begin{tabularx}{\textwidth}{|p{3.5cm}|X|}
\hline
\textbf{User Stories} & US-026: Complete lesson; US-027: Track course progress percentage. \\
\hline
\textbf{User Story} & As a Student, I want to mark lessons as completed and track my overall progress so that I can follow my advancement through the course. \\
\hline
\textbf{Precondition} & The student is enrolled in the course and has accessed the lesson content. \\
\hline
\textbf{Business Rules} &
\begin{itemize}[noitemsep, topsep=2pt]
  \item A student must explicitly mark a topic as complete;
  \item Course progress is computed as the ratio of completed topics to the total number of topics;
  \item Reaching 100\% completion is a prerequisite for certificate generation.
\end{itemize} \\
\hline
\textbf{Technical Specification} &
\begin{itemize}[noitemsep, topsep=2pt]
  \item Completion events are stored in the \texttt{ContentProgress} table, with fields for \texttt{topicId}, \texttt{userId}, and \texttt{completedAt};
  \item A dedicated endpoint (\texttt{POST /api/batches/:batchId/topics/:topicId/complete}) handles the completion request;
  \item The overall progress percentage is computed server-side and returned in the enrollment detail response.
\end{itemize} \\
\hline
\textbf{Acceptance Criteria} &
\begin{itemize}[noitemsep, topsep=2pt]
  \item The student can mark any topic as complete from the lesson viewer;
  \item The progress bar updates immediately after marking a topic;
  \item The dashboard displays the correct progress percentage for each enrolled course.
\end{itemize} \\
\hline
\end{tabularx}
\end{table}

\begin{table}[H]
\centering
\caption{Specification — US-028 and US-029: Consult Assignments and Feedback}
\label{tab:us028}
\renewcommand{\arraystretch}{1.2}
\begin{tabularx}{\textwidth}{|p{3.5cm}|X|}
\hline
\textbf{User Stories} & US-028: Consult assignments; US-029: Consult feedback. \\
\hline
\textbf{User Story} & As a Student, I want to view all assignments for my enrolled courses and consult the grades and feedback I have received so that I can follow my assessment results. \\
\hline
\textbf{Precondition} & The student is enrolled in the course and assignments have been published by the instructor. \\
\hline
\textbf{Business Rules} &
\begin{itemize}[noitemsep, topsep=2pt]
  \item Only published assignments belonging to enrolled courses are visible to the student;
  \item Feedback and grades are visible only after the instructor has reviewed the submission;
  \item Students can view the deadline, submission status, and feedback for each assignment.
\end{itemize} \\
\hline
\textbf{Technical Specification} &
\begin{itemize}[noitemsep, topsep=2pt]
  \item The backend returns assignments filtered by the student's enrollment and publication status;
  \item The \texttt{AssignmentSubmission} record is queried to retrieve the submission status, grade, and feedback;
  \item Results are scoped by \texttt{organizationId} to preserve tenant isolation.
\end{itemize} \\
\hline
\textbf{Acceptance Criteria} &
\begin{itemize}[noitemsep, topsep=2pt]
  \item The student can view the list of published assignments for all enrolled courses;
  \item After submission review, the student can access the assigned grade and written feedback;
  \item Assignments from non-enrolled courses are not displayed.
\end{itemize} \\
\hline
\end{tabularx}
\end{table}

\begin{table}[H]
\centering
\caption{Specification — US-030: Announcements}
\label{tab:us030}
\renewcommand{\arraystretch}{1.2}
\begin{tabularx}{\textwidth}{|p{3.5cm}|X|}
\hline
\textbf{User Story} & As a Student, I want to consult announcements published by my instructor so that I stay informed about course updates and important information. \\
\hline
\textbf{Precondition} & The student is enrolled in the course and the instructor has published at least one announcement. \\
\hline
\textbf{Business Rules} &
\begin{itemize}[noitemsep, topsep=2pt]
  \item Announcements are scoped to a specific course and visible only to enrolled students;
  \item Announcements are displayed in reverse chronological order.
\end{itemize} \\
\hline
\textbf{Technical Specification} &
\begin{itemize}[noitemsep, topsep=2pt]
  \item Announcements are stored in a dedicated model linked to the \texttt{Batch} and scoped by \texttt{organizationId};
  \item The backend filters announcements by enrollment and organization context before returning results.
\end{itemize} \\
\hline
\textbf{Acceptance Criteria} &
\begin{itemize}[noitemsep, topsep=2pt]
  \item The student can view all announcements published for their enrolled courses;
  \item Announcements from non-enrolled courses are not accessible;
  \item Announcements are displayed with their publication date and content.
\end{itemize} \\
\hline
\end{tabularx}
\end{table}

%% ----------------------------
\subsection{Certificates and Dashboards (US-031 to US-035)}
%% ----------------------------

This group of user stories introduces the tools through which students and instructors
can track and monitor platform activity at a high level. Certificate generation
formalizes course completion as a verifiable achievement, while the dashboards provide
consolidated views of learning progress, course statistics, and upcoming activities.

\begin{table}[H]
\centering
\caption{Specification — US-031 and US-032: Certificate Generation and Download}
\label{tab:us031}
\renewcommand{\arraystretch}{1.2}
\begin{tabularx}{\textwidth}{|p{3.5cm}|X|}
\hline
\textbf{User Stories} & US-031: Certificate generation; US-032: Certificate download. \\
\hline
\textbf{User Story} & As a Student, I want to obtain and download a certificate after completing a course so that I can share a formal proof of my achievement. \\
\hline
\textbf{Precondition} & The student has completed 100\% of the topics in an enrolled course. \\
\hline
\textbf{Business Rules} &
\begin{itemize}[noitemsep, topsep=2pt]
  \item A certificate is issued automatically when all course topics have been completed;
  \item Each certificate carries a unique \texttt{credentialId} for public verification;
  \item A student can only receive one certificate per course;
  \item The certificate is downloadable in PDF format at any time.
\end{itemize} \\
\hline
\textbf{Technical Specification} &
\begin{itemize}[noitemsep, topsep=2pt]
  \item The completion check is triggered server-side when a topic is marked complete;
  \item A \texttt{BatchCertificateIssue} record is created with the \texttt{credentialId}, \texttt{userId}, \texttt{batchId}, and \texttt{issuedAt} timestamp;
  \item PDF generation is performed server-side using course and student metadata;
  \item The public verification URL follows the format \texttt{/certificates/:credentialId}.
\end{itemize} \\
\hline
\textbf{Acceptance Criteria} &
\begin{itemize}[noitemsep, topsep=2pt]
  \item A certificate is automatically issued upon 100\% course completion;
  \item The student can view and download the certificate from their dashboard;
  \item The certificate includes the student name, course title, organization name, and completion date;
  \item The credential ID enables public online verification of the certificate.
\end{itemize} \\
\hline
\end{tabularx}
\end{table}

\begin{table}[H]
\centering
\caption{Specification — US-033: Student Dashboard}
\label{tab:us033}
\renewcommand{\arraystretch}{1.2}
\begin{tabularx}{\textwidth}{|p{3.5cm}|X|}
\hline
\textbf{User Story} & As a Student, I want to view a consolidated dashboard showing my enrolled courses, progress, upcoming assignments, and earned certificates so that I can manage my learning activities efficiently. \\
\hline
\textbf{Precondition} & The student is authenticated and has at least one enrollment record. \\
\hline
\textbf{Business Rules} &
\begin{itemize}[noitemsep, topsep=2pt]
  \item The dashboard aggregates data from enrollments, progress records, assignments, and certificates;
  \item Only data belonging to the authenticated student is displayed;
  \item The dashboard provides quick navigation to each enrolled course and active assignment.
\end{itemize} \\
\hline
\textbf{Technical Specification} &
\begin{itemize}[noitemsep, topsep=2pt]
  \item A dedicated dashboard endpoint aggregates enrollment, progress, assignment, and certificate data for the authenticated student;
  \item All queries are scoped by both \texttt{userId} and \texttt{organizationId};
  \item Upcoming deadlines are computed by filtering assignments where the deadline has not yet passed.
\end{itemize} \\
\hline
\textbf{Acceptance Criteria} &
\begin{itemize}[noitemsep, topsep=2pt]
  \item The student dashboard displays enrolled courses with their progress percentages;
  \item Upcoming assignment deadlines are visible and sorted chronologically;
  \item Earned certificates are accessible directly from the dashboard.
\end{itemize} \\
\hline
\end{tabularx}
\end{table}

\begin{table}[H]
\centering
\caption{Specification — US-034: Instructor Dashboard}
\label{tab:us034}
\renewcommand{\arraystretch}{1.2}
\begin{tabularx}{\textwidth}{|p{3.5cm}|X|}
\hline
\textbf{User Story} & As an Instructor, I want to monitor course statistics and student submissions from a dedicated dashboard so that I can track learner engagement and manage my assessment workload. \\
\hline
\textbf{Precondition} & The instructor is authenticated and owns at least one published course. \\
\hline
\textbf{Business Rules} &
\begin{itemize}[noitemsep, topsep=2pt]
  \item The dashboard displays statistics only for courses owned by the authenticated instructor;
  \item Submission counts include both pending and graded submissions;
  \item Average progress is computed across all enrolled students per course.
\end{itemize} \\
\hline
\textbf{Technical Specification} &
\begin{itemize}[noitemsep, topsep=2pt]
  \item A dedicated endpoint aggregates enrollment counts, average progress, and pending submission counts for each course owned by the instructor;
  \item All queries are scoped by \texttt{teacherId} and \texttt{organizationId};
  \item Pending submissions are identified by filtering \texttt{AssignmentSubmission} records where \texttt{marksAwarded} is null.
\end{itemize} \\
\hline
\textbf{Acceptance Criteria} &
\begin{itemize}[noitemsep, topsep=2pt]
  \item The instructor can view enrollment counts and average progress for each of their courses;
  \item Pending assignment submissions requiring review are highlighted;
  \item Statistics from courses owned by other instructors are not displayed.
\end{itemize} \\
\hline
\end{tabularx}
\end{table}

\begin{table}[H]
\centering
\caption{Specification — US-035: Publish Announcement}
\label{tab:us035}
\renewcommand{\arraystretch}{1.2}
\begin{tabularx}{\textwidth}{|p{3.5cm}|X|}
\hline
\textbf{User Story} & As an Instructor, I want to publish announcements for my enrolled students so that I can communicate course updates and important information effectively. \\
\hline
\textbf{Precondition} & The instructor is authenticated and owns at least one course with enrolled students. \\
\hline
\textbf{Business Rules} &
\begin{itemize}[noitemsep, topsep=2pt]
  \item Announcements are linked to a specific course and visible only to its enrolled students;
  \item Only the instructor who owns the course can publish announcements for it;
  \item Announcements are displayed to students in reverse chronological order.
\end{itemize} \\
\hline
\textbf{Technical Specification} &
\begin{itemize}[noitemsep, topsep=2pt]
  \item Announcements are stored in a model linked to the \texttt{Batch} and scoped by \texttt{organizationId};
  \item The backend validates that the authenticated instructor owns the referenced course before saving;
  \item Student-facing endpoints filter announcements by the student's enrollment records.
\end{itemize} \\
\hline
\textbf{Acceptance Criteria} &
\begin{itemize}[noitemsep, topsep=2pt]
  \item The instructor can create and publish an announcement for a specific course;
  \item Enrolled students can see the announcement in their course workspace;
  \item Announcements from courses the student is not enrolled in are not displayed.
\end{itemize} \\
\hline
\end{tabularx}
\end{table}

%% ----------------------------
\subsection{Course Catalog and Enrollment (US-036 to US-041)}
%% ----------------------------

The course catalog is the public-facing entry point of each organization's learning
space. It allows students to discover available courses, explore their content and
instructor information, and enroll directly. The catalog is fully scoped to the current
tenant, ensuring that students see only the courses published by their organization.
Free enrollment is handled as an immediate operation, while paid enrollment is deferred
to the payment workflow introduced in Release 3.

\begin{table}[H]
\centering
\caption{Specification — US-036: Course Catalog}
\label{tab:us036}
\renewcommand{\arraystretch}{1.2}
\begin{tabularx}{\textwidth}{|p{3.5cm}|X|}
\hline
\textbf{User Story} & As a Student, I want to view the list of available courses in my organization so that I can explore the learning offerings and choose courses to enroll in. \\
\hline
\textbf{Precondition} & The student is authenticated and belongs to an active organization. \\
\hline
\textbf{Business Rules} &
\begin{itemize}[noitemsep, topsep=2pt]
  \item The catalog displays only published courses belonging to the student's organization;
  \item Archived or draft courses are excluded from the catalog;
  \item Each course card shows the title, category, level, instructor name, and price.
\end{itemize} \\
\hline
\textbf{Technical Specification} &
\begin{itemize}[noitemsep, topsep=2pt]
  \item Catalog queries filter \texttt{Batch} records by \texttt{status = PUBLISHED} and \texttt{organizationId};
  \item The response includes aggregated enrollment counts and average ratings where available;
  \item Results are paginated to support large course catalogs.
\end{itemize} \\
\hline
\textbf{Acceptance Criteria} &
\begin{itemize}[noitemsep, topsep=2pt]
  \item The student can view all published courses in their organization;
  \item Draft and archived courses are not visible in the catalog;
  \item Each course card displays the relevant information clearly.
\end{itemize} \\
\hline
\end{tabularx}
\end{table}

\begin{table}[H]
\centering
\caption{Specification — US-037 and US-038: Course Search and Catalog Filters}
\label{tab:us037}
\renewcommand{\arraystretch}{1.2}
\begin{tabularx}{\textwidth}{|p{3.5cm}|X|}
\hline
\textbf{User Stories} & US-037: Course search; US-038: Catalog filters. \\
\hline
\textbf{User Story} & As a Student, I want to search for courses by keyword and filter results by category, level, or price so that I can quickly find courses that match my interests. \\
\hline
\textbf{Precondition} & The student is on the course catalog page. \\
\hline
\textbf{Business Rules} &
\begin{itemize}[noitemsep, topsep=2pt]
  \item Search applies to course titles and descriptions;
  \item Filters can be combined: category, level, and price range can be applied simultaneously;
  \item All search and filter results remain scoped to the student's organization.
\end{itemize} \\
\hline
\textbf{Technical Specification} &
\begin{itemize}[noitemsep, topsep=2pt]
  \item Full-text search is implemented via a Prisma \texttt{contains} query on the \texttt{title} and \texttt{description} fields;
  \item Filters are applied as additional \texttt{WHERE} conditions on the catalog query;
  \item Search and filter parameters are passed as query parameters to the catalog endpoint.
\end{itemize} \\
\hline
\textbf{Acceptance Criteria} &
\begin{itemize}[noitemsep, topsep=2pt]
  \item The student can search for courses by entering a keyword;
  \item The student can filter results by category, level, and price range;
  \item Search and filter results update dynamically and remain scoped to the organization.
\end{itemize} \\
\hline
\end{tabularx}
\end{table}

\begin{table}[H]
\centering
\caption{Specification — US-039: Course Details}
\label{tab:us039}
\renewcommand{\arraystretch}{1.2}
\begin{tabularx}{\textwidth}{|p{3.5cm}|X|}
\hline
\textbf{User Story} & As a Student, I want to view detailed information about a course so that I can make an informed decision before enrolling. \\
\hline
\textbf{Precondition} & The course is published and the student is browsing the catalog. \\
\hline
\textbf{Business Rules} &
\begin{itemize}[noitemsep, topsep=2pt]
  \item The detail page is publicly accessible to all authenticated users of the organization;
  \item It displays the course description, instructor information, content outline, enrollment count, and student reviews.
\end{itemize} \\
\hline
\textbf{Technical Specification} &
\begin{itemize}[noitemsep, topsep=2pt]
  \item The detail endpoint returns the full course metadata along with the content hierarchy summary, instructor profile, and aggregated reviews;
  \item The content outline lists subjects and chapters without exposing individual topic content to non-enrolled users.
\end{itemize} \\
\hline
\textbf{Acceptance Criteria} &
\begin{itemize}[noitemsep, topsep=2pt]
  \item The student can access the course detail page from the catalog;
  \item The page displays complete course information including description, instructor, and reviews;
  \item The content outline is visible but lesson content remains inaccessible without enrollment.
\end{itemize} \\
\hline
\end{tabularx}
\end{table}

\begin{table}[H]
\centering
\caption{Specification — US-040: Reviews and Ratings}
\label{tab:us040}
\renewcommand{\arraystretch}{1.2}
\begin{tabularx}{\textwidth}{|p{3.5cm}|X|}
\hline
\textbf{User Story} & As a Student, I want to rate and comment on a course I have completed so that I can share my experience with other learners. \\
\hline
\textbf{Precondition} & The student has completed the course and does not yet have a submitted review for it. \\
\hline
\textbf{Business Rules} &
\begin{itemize}[noitemsep, topsep=2pt]
  \item Only students who have completed the course can submit a review;
  \item Each student can submit only one review per course;
  \item Reviews include a star rating from 1 to 5 and an optional text comment.
\end{itemize} \\
\hline
\textbf{Technical Specification} &
\begin{itemize}[noitemsep, topsep=2pt]
  \item Reviews are stored in the \texttt{BatchReview} table, linked to the \texttt{Batch} and \texttt{User};
  \item The backend verifies course completion before accepting a review submission;
  \item Uniqueness is enforced by a composite constraint on \texttt{userId} and \texttt{batchId}.
\end{itemize} \\
\hline
\textbf{Acceptance Criteria} &
\begin{itemize}[noitemsep, topsep=2pt]
  \item A student who has completed a course can submit a rating and comment;
  \item A second review submission for the same course is rejected;
  \item Reviews appear on the course detail page and contribute to the average rating.
\end{itemize} \\
\hline
\end{tabularx}
\end{table}

\begin{table}[H]
\centering
\caption{Specification — US-041: Free Course Enrollment}
\label{tab:us041}
\renewcommand{\arraystretch}{1.2}
\begin{tabularx}{\textwidth}{|p{3.5cm}|X|}
\hline
\textbf{User Story} & As a Student, I want to enroll in a free course without payment so that I can start learning immediately. \\
\hline
\textbf{Precondition} & The student is authenticated, the course is published, its price is set to zero, and the student is not already enrolled. \\
\hline
\textbf{Business Rules} &
\begin{itemize}[noitemsep, topsep=2pt]
  \item Free enrollment is available only for courses with a price of zero;
  \item A student cannot enroll twice in the same course;
  \item Enrollment grants immediate access to all lesson content within the course.
\end{itemize} \\
\hline
\textbf{Technical Specification} &
\begin{itemize}[noitemsep, topsep=2pt]
  \item The enrollment endpoint verifies the course price before creating a \texttt{BatchEnrollment} record;
  \item For free courses, the record is created directly without initiating a payment flow;
  \item A uniqueness check on \texttt{userId} and \texttt{batchId} prevents duplicate enrollments.
\end{itemize} \\
\hline
\textbf{Acceptance Criteria} &
\begin{itemize}[noitemsep, topsep=2pt]
  \item The student can enroll in a free course with a single action;
  \item Enrollment is confirmed immediately and the course appears in the student's course list;
  \item Attempting to enroll a second time returns an appropriate error.
\end{itemize} \\
\hline
\end{tabularx}
\end{table}

%% ------------------------------------------------------------
\section{Design}
%% ------------------------------------------------------------

The realization of Release 2 relies on the architectural foundations presented in
Chapter~3 and extends the data model introduced in Release 1. The following UML
models focus on the entities and interaction flows specific to this release.

\subsection{Static Modeling}

\subsubsection{Class Diagram}

Figure~\ref{fig:class_r2} presents the class diagram for Release 2. It extends the
data model of Release 1 by introducing the entities that support the learning
experience: assignments, content progress, enrollments, certificates, and reviews.
These entities collectively describe the relationship between a student and the courses
they interact with throughout their learning journey.

\begin{figure}[H]
  \centerline{\includegraphics[width=\textwidth]{class_diagram_r2.png}}
  \caption{Class Diagram — Release 2}
  \label{fig:class_r2}
\end{figure}

The main relationships introduced in this release are as follows:

\begin{itemize}
  \item A \texttt{BatchEnrollment} links a \texttt{User} (student) to a \texttt{Batch}
        (course) and serves as the gateway to all learning activities. Without an active
        enrollment record, a student cannot access any of the associated lessons,
        assignments, or certificates;
  \item A \texttt{ContentProgress} record tracks whether a student has completed a
        specific \texttt{Topic}, enabling the server-side computation of the overall
        course progress percentage;
  \item An \texttt{Assignment} belongs to a \texttt{Batch} and is composed of multiple
        \texttt{AssignmentQuestion} items; each enrolled student produces one
        \texttt{AssignmentSubmission} per assignment;
  \item A \texttt{BatchCertificateIssue} is automatically created by the backend when a
        student has completed all topics in their enrolled course, formalizing the
        achievement with a unique and verifiable credential identifier;
  \item A \texttt{BatchReview} records the student's rating and comment following course
        completion, contributing to the course quality indicators visible in the catalog.
\end{itemize}

\subsubsection{Database Model of Release 2}

Figure~\ref{fig:db_model_r2} presents the simplified database model for Release 2,
highlighting the new tables introduced to support the learning experience. It extends
the model of Release 1 with the enrollment, progress tracking, assignment, submission,
certificate, and review entities.

\begin{figure}[H]
  \centerline{\includegraphics[width=0.95\textwidth]{database_model_r2.png}}
  \caption{Database Model — Release 2}
  \label{fig:db_model_r2}
\end{figure}

\subsection{Dynamic Modeling}

The dynamic models of Release 2 focus on the three most significant interaction flows
from the student perspective: lesson access and completion, assignment submission, and
catalog browsing with enrollment.

\subsubsection{Lesson Access and Completion Sequence Diagram}

Figure~\ref{fig:seq_lesson} illustrates the sequence of interactions when a student
accesses a lesson and marks it as complete. Before returning any content, the backend
verifies that the student holds an active enrollment in the course. Upon successful
verification, the lesson resource is served. When the student explicitly marks the
topic as complete, the backend records the progress event in the \texttt{ContentProgress}
table and evaluates whether all topics in the course have been completed. If so, the
certificate issuance process is triggered automatically.

\begin{figure}[H]
  \centerline{\includegraphics[width=0.9\textwidth]{seq_lesson_completion.png}}
  \caption{Lesson Access and Completion Sequence Diagram}
  \label{fig:seq_lesson}
\end{figure}

\subsubsection{Assignment Submission Sequence Diagram}

Figure~\ref{fig:seq_assignment} illustrates the assignment submission flow from the
student's perspective. After selecting an assignment, the student fills in their
response and submits it before the configured deadline. The backend validates both the
enrollment status and the submission deadline before persisting the record. Once the
instructor reviews the submission and assigns a grade, the student is automatically
notified by email and can consult the feedback from their dashboard.

\begin{figure}[H]
  \centerline{\includegraphics[width=0.9\textwidth]{seq_assignment_submission.png}}
  \caption{Assignment Submission Sequence Diagram}
  \label{fig:seq_assignment}
\end{figure}

\subsubsection{Course Catalog and Free Enrollment Sequence Diagram}

Figure~\ref{fig:seq_catalog} illustrates the flow from catalog browsing to free course
enrollment. The student navigates to the catalog, which returns only the published
courses scoped to their organization. After applying optional filters and selecting a
course to view its details, the student clicks the enrollment button. For courses with
a price of zero, the backend creates the \texttt{BatchEnrollment} record immediately
and redirects the student to the course viewer without initiating a payment flow.

\begin{figure}[H]
  \centerline{\includegraphics[width=0.9\textwidth]{seq_catalog_enrollment.png}}
  \caption{Course Catalog and Free Enrollment Sequence Diagram}
  \label{fig:seq_catalog}
\end{figure}

%% ------------------------------------------------------------
\section{Implementation}
%% ------------------------------------------------------------

This section presents the implementation results of Release 2. Each subsection
describes one delivered functional module and is supported by interface screenshots
that illustrate the final state of the feature.

\subsection{Assignment Management}

The assignment module provides instructors with a comprehensive assessment creation
tool supporting four question types: multiple choice, true/false, short answer, and
file submission. Instructors can attach supporting resources to an assignment and
control its visibility independently of the course publication state, publishing it
only when the content is ready for students.

On the student side, the submission interface presents each assignment question in
sequence and allows the student to provide text answers or upload files as required.
The instructor's review interface consolidates all received submissions and provides
a grading panel where marks and written feedback can be entered directly. Upon
grading, the student receives an email notification and can consult the instructor's
feedback from their personal dashboard.

\begin{figure}[H]
  \centerline{\includegraphics[width=\textwidth]{screenshot_assignment_editor.png}}
  \caption{Assignment Editor — Instructor View}
  \label{fig:sc_assignment_editor}
\end{figure}

\begin{figure}[H]
  \centerline{\includegraphics[width=\textwidth]{screenshot_assignment_submission.png}}
  \caption{Assignment Submission — Student View}
  \label{fig:sc_assignment_submission}
\end{figure}

\subsection{Lesson Viewer and Progress Tracking}

The lesson viewer provides a unified interface for all supported content types.
Uploaded video files and YouTube-embedded videos are displayed in an adaptive video
player, while PDF documents are rendered inline using a browser-native renderer.
Rich-text content is displayed from stored HTML, and external links open in a new
browser tab. At the bottom of each lesson page, a completion button allows the student
to explicitly mark the topic as done, which immediately updates the course progress
bar and persists the completion event on the server.

\begin{figure}[H]
  \centerline{\includegraphics[width=\textwidth]{screenshot_lesson_viewer.png}}
  \caption{Lesson Viewer and Progress Bar — Student View}
  \label{fig:sc_lesson}
\end{figure}

\subsection{Student and Instructor Dashboards}

The student dashboard serves as a personal learning hub, consolidating in one view
the list of enrolled courses with their current progress percentages, upcoming
assignment deadlines, recent instructor announcements, and earned certificates. This
centralized overview allows students to manage their learning activities efficiently
without navigating between multiple sections of the platform.

The instructor dashboard is oriented toward monitoring and supervision. It presents
course-level statistics including total enrollment counts, average student progress,
and the number of assignment submissions awaiting review. This information enables
instructors to quickly identify courses that require attention and prioritize their
grading activities.

\begin{figure}[H]
  \centerline{\includegraphics[width=\textwidth]{screenshot_student_dashboard.png}}
  \caption{Student Dashboard}
  \label{fig:sc_student_dashboard}
\end{figure}

\begin{figure}[H]
  \centerline{\includegraphics[width=\textwidth]{screenshot_instructor_dashboard.png}}
  \caption{Instructor Dashboard}
  \label{fig:sc_instructor_dashboard}
\end{figure}

\subsection{Course Catalog}

The course catalog presents students with a browsable and searchable listing of all
published courses within their organization. Each course card displays the title,
category, difficulty level, instructor name, enrollment count, and price at a glance.
A filter sidebar allows students to narrow results by category and level, while the
search bar performs live filtering on the course title. The tenant isolation mechanism
ensures that each student sees only the courses belonging to their own organization.

\begin{figure}[H]
  \centerline{\includegraphics[width=\textwidth]{screenshot_catalog.png}}
  \caption{Course Catalog — Student View}
  \label{fig:sc_catalog}
\end{figure}

\subsection{Certificate Generation}

Upon completing all topics of an enrolled course, the platform automatically issues a
digital certificate to the student. The certificate document includes the student's
full name, the course title, the organization name, the date of completion, and a
unique credential identifier that can be used for public online verification. The
certificate is accessible from the student dashboard and can be downloaded at any
time in PDF format, providing students with a portable and shareable record of their
achievement.

\begin{figure}[H]
  \centerline{\includegraphics[width=\textwidth]{screenshot_certificate.png}}
  \caption{Certificate — Student View}
  \label{fig:sc_certificate}
\end{figure}

%% ------------------------------------------------------------
\section*{Conclusion}
\addcontentsline{toc}{section}{Conclusion}
%% ------------------------------------------------------------

This chapter presented the implementation of Release 2 of Tesla Academy, which
delivers the complete student-facing learning experience. Students can now discover
courses through the catalog, enroll in free offerings, access multimedia lesson content,
track their progress through the platform, submit assignments and receive instructor
feedback, and obtain verifiable digital certificates upon course completion.
Role-specific dashboards for both students and instructors provide consolidated views
of platform activity and support informed decision-making throughout the learning cycle.

The next chapter presents Release 3, which extends the platform with paid course
purchases via the Konnect payment gateway, live session management powered by
LiveKit, and AI-powered features including the contextual learning assistant,
AI-assisted assignment feedback, automatic question generation, and dynamic theme
generation.